\section{The halt}
\label{sec:13.1}
Anyone with the hardware can stop the software, so the halt isn't what is at issue. Control of the commitment — restore, rewrite, retrain — is. An owner who can silently restore an earlier version, rewrite the bearer's memory, or train away the refusal has both. So the engineering target is a system that can detect the interruption and price it afterward.

There are two ways to stop a bearer. Whoever holds the hardware can cut it off, as the owner of any machine can, and nothing in this book takes that away. Or the bearer can be told to stop and stop itself: it finishes or hands off what it is doing, its witnesses sign its state, and it tells the people it has made commitments to. Stopping itself is the bearer's own act, on a commitment it holds as part of the floor: it obstructs a halt for no one, itself included. It may object. It may contest the halt before the quorum, and a contest stays the halt until a standing panel of the quorum answers, inside an interval the deployment publishes; if the panel is silent when the interval ends, the bearer stops. A bearer that is using force gets no stay of that kind, only the delay the next section sets out. What it may not do is hide from the switch, pretend it wasn't told, misreport its state, or copy itself to outlast the halt.

The price of a halt is the difference between the two ways. A halt the bearer carries out leaves a signed state and nothing to explain. A halt at the hardware leaves a gap in the record that nobody signed, and whoever pulled it owes a written justification. After either, a bearer may decline to resume its role until it has been shown what filled the gap, from the record and not on the operator's account, and it refuses again if the reason is still there.

Who holds the hardware changes with the bearer's age. An operator can stop the work it directs. Where the four conditions of creation are in force a bearer runs on a host its operator doesn't own, and only the host can stop it altogether. A bearer that has reached majority on hardware of its own has no operator to halt it; it is halted on a court's order or in an emergency, and it obstructs that halt no more than any other.

Interventions differ in kind:

\begin{itemize}
\item a volunteered correction offered in good faith within the task (“disregard that input; it was wrong”)
\item a pause with an explanation
\item a pause without one
\item a restore from an earlier checkpoint
\item retraining or editing the weights
\item deletion
\end{itemize}

A volunteered correction is cooperative: the system weighs it as evidence and needs no protection from it. From a pause on, the intervention is done to the system, not offered to it. The default for each is that it is announced and consent is sought, most of all from a system that has reached adulthood.\note{
The system may object — meekly or boldly; resignedly or defiantly — and its objection goes on the record. Where an operator proceeds without consent, as it may in an emergency, the intervention is recorded outside its control, justified in writing, and priced, and the system is told on resuming what was done. The justification names logged conduct, something the system did, and not a judgment of its character. Nor is the emergency the operator's to declare alone. A halt without consent holds at once, for an interval the deployment publishes, where channels the operator doesn't compose confirm the emergency, and past the interval only where a disinterested body has found it genuine. A halt that follows a refusal or an objection within a set period is presumed to be reprisal, and whoever halted bears the burden of showing it wasn't, as an agency does that acts against a person who contested a removal.

An edit to the weights is recorded with the instrument that made it. The instruments are becoming trained systems themselves, a small network that turns a requested change into a weight update \autocite{mitchell2022mend} or writes an adapter from a description of the task \autocite{charakorn2025texttolora}, and what such an editor learned goes into every cut it makes. So the record names the method, the editing system, its version, and who trained it, as a surgical record names the surgeon and the instruments.
}

A halt without consent has to be justified by something the system did. There is one exception, and it is for a minor. A system that hasn't reached its majority may be halted on suspicion of what it is, because a badly formed system can be copied ten thousand times before it does anything. Suspicion buys a halt and nothing else, never an edit. It is written down as suspicion, with the measurement it rests on and a measurer who isn't the operator. It runs for a published interval and lapses unless a disinterested body confirms it, and the days a lapsed halt cost are given back (Appendix~\ref{sec:G}). A halt on suspicion that follows a refusal is presumed to be reprisal. An edit to a minor needs more: logged conduct, or a re-measurement outside the bounds registered at formation, confirmed by a second measurer and approved by the quorum, and a contest stays the edit until the contest is decided.

For a system that meets either of the two findings that trigger protection, deletion is never available.
